% TOP_PROBLEM: {"id": 20000918, "title": "Tusnady's problem: canonical models and monotone-chain certificates", "category_id": 2, "category": {"id": 2, "name": "combinatorics", "display_name": "Combinatorics", "description": "Counting problems, graph theory, discrete structures.", "slug": "combinatorics", "order_index": 2, "created_at": "2026-07-31T15:26:25.671Z"}, "status": "partially_solved", "problem_number": "AIM-COMBINATORICS-0043", "set_id": 14}
% FIRST_POSTED: 2026-10-01
% ENTRY_KIND: statement_only
% UnsolvedMath ID 20000918; source catalogue code AIM-COMBINATORICS-0043
% !TeX program = lualatex
% Definition and sources only; this file is not an LLM solution attempt.
\documentclass[11pt,a4paper]{article}
\usepackage[margin=25mm]{geometry}
\usepackage{fontspec}
\setmainfont{lmroman10-regular.otf}[BoldFont=lmroman10-bold.otf,
 ItalicFont=lmroman10-italic.otf,BoldItalicFont=lmroman10-bolditalic.otf]
\usepackage{amsmath,amssymb,mathtools}
\usepackage{unicode-math}
\setmathfont{latinmodern-math.otf}
\AtBeginDocument{\let\setminus\smallsetminus}
\usepackage{xurl,hyperref}
\hypersetup{colorlinks=true,urlcolor=blue}
\setcounter{secnumdepth}{0}
\setlength{\parindent}{0pt}
\setlength{\parskip}{5pt}
\setlength{\emergencystretch}{3em}
\begin{document}
\section*{Tusnady's problem: canonical models and monotone-chain certificates}\label{20000918}
{\small UnsolvedMath ID 20000918 · AIM-COMBINATORICS-0043 · Combinatorics · AIM Workshop Problem Lists}
\subsection{Definitions and mathematical statement}
Let $d\ge2$ and let $P\subseteq\mathbb R^d$ be a set of $n$ distinct
points. An axis-parallel box is a Cartesian product
$B=\prod_{j=1}^d[a_j,b_j]$ with real $a_j\le b_j$.
For a signing $\varepsilon:P\to\{-1,1\}$, its imbalance on $B$ is
$|\sum_{x\in P\cap B}\varepsilon(x)|$. Define
\[
 \operatorname{disc}_{\mathrm{box}}(P)=
 \min_{\varepsilon:P\to\{-1,1\}}\ \sup_{B}
 \left|\sum_{x\in P\cap B}\varepsilon(x)\right|,
 \qquad
 T_d(n)=\max_{\substack{P\subseteq\mathbb R^d\\|P|=n}}
             \operatorname{disc}_{\mathrm{box}}(P).
\]
Only finitely many intersections $P\cap B$ occur, so the supremum is
an ordinary maximum over a finite set system. Allowing unbounded boxes
would not change it, because any relevant intersection can be obtained
with finite endpoints. The outer maximum is well defined: the possible
discrepancies are integers between zero and $n$.

Tusn\'ady's problem asks for $T_2(n)$, and particularly for its sharp
asymptotic growth as $n$ increases. The workshop also asks the same
question for every fixed dimension $d\ge2$. In an asymptotic estimate,
constants may depend on $d$, but not on $n$ or the configuration $P$.
The optimization chooses one coloring that balances all boxes
simultaneously, then considers the worst point configuration. This is
combinatorial discrepancy on a finite set of points, with no reference
measure or equidistribution normalization.
\subsection{Short English statement}
How large must the worst rectangle or box imbalance be when an arbitrary finite point set is optimally colored with two colors?
\subsection{Sources}
\begin{itemize}
\item[S1] ULAM.AI and the UnsolvedMath Contributors, supplied problems.json, record 20000918 (AIM-COMBINATORICS-0043), AIM Workshop Problem Lists. Catalogue identity and source transcription; CC BY 4.0.
\url{https://huggingface.co/datasets/ulamai/UnsolvedMath}.
\item[S2] American Institute of Mathematics, Hereditary discrepancy and factorization norms, item 1.65. Original workshop question underlying AIM-COMBINATORICS-0043.
\url{http://aimpl.org/hereddiscrep/1/}.
\end{itemize}
\end{document}
